\section{Conclusion}

We presented a falsification-first, galaxy-domain response-sector framework that makes an ``alternative to particulate CDM'' precise within standard GR. The response sector \textbf{(see PDF; response-sector symbol not recovered from DOCX extraction)} is defined as the effective stress--energy required so that the metric potentials inferred from galaxy phenomenology satisfy \textbf{(see PDF; Einstein equation not recovered from DOCX extraction)} with baryons fixed by SPARC mass models. As emphasized in \S9, this definition is ontologically agnostic and is best read as an EFT/constitutive-closure candidate rather than a literal mechanical medium.

Using SPARC-175, we showed that a low-DOF boundary-layer closure producing an asymptotic \textbf{(see PDF; tail notation not recovered)} tail yields strong improvement over baryons-only fits across the sample and recovers the BTFR as an empirical consequence. Robustness is strengthened by separating \textbf{(see PDF; model-fitted amplitude symbol not recovered)} from a model-free outer estimator \textbf{(see PDF; outer estimator symbol not recovered)}; their strong agreement in rank supports that the inferred amplitude is not an artifact of inner weighting, while discrepancies diagnose inner--outer shape tension and motivate minimal extension (e.g., a width parameter).

The work is deliberately scoped. Several limitations stated in \S9 remain open: the activation/closure is phenomenological rather than first-principles; lensing consistency is a working hypothesis pending measurement; clusters and pressure-supported systems require extension; and cosmological viability is an explicit requirement rather than a demonstrated result. These are not rhetorical caveats but the defined research program.

Accordingly, we pre-register a validation ladder:
\begin{itemize}
	\item continued compression and cross-validation on rotation curves;
	\item lensing adjudication of Branch A (\textbf{see PDF; metric-equivalent condition}) versus Branch B (slip-allowed, \textbf{see PDF; slip condition} in activation zones);
	\item cluster-regime scaling tests; and
	\item perturbation-level cosmological consistency.
\end{itemize}
